\documentclass[11pt]{ctexart}
\usepackage[a4paper,margin=1in]{geometry}
\usepackage{graphicx}
\usepackage{xcolor}
\usepackage{hyperref}
\definecolor{refgray}{gray}{0.45}
\title{\textbf{市场最新观点汇总}\\[0.4em]{\large 全球资产定价正从地缘风险折价转向结构性再定价，AI驱动的盈利扩散与供应链重构成为主线。}}
\author{KC桌面 自动生成}
\date{260516}
\begin{document}
\maketitle
\section*{一页摘要}
\begin{itemize}
  \item 美股盈利增长正从Mag 7向S\&P 493扩散，盈利复苏的广度被市场低估。
  \item 中国经常账户顺差韧性超预期，服务逆差收窄成为新的支撑变量。
  \item 人形机器人出海加速，中国供应链的全球扩张成为产业主线。
  \item 存储芯片正经历从周期性到成长性的转变，AI推理需求带来指数级增长。
  \item 人民币资产的风险溢价正在重构，贸易与投资实质性松动提供支撑。
  \item 新兴市场外汇面临结构性压力，印尼盾等货币脆弱性突出。
  \item 美国通胀正从临时因素转变为结构性成本传递，美联储政策路径可能更鹰。
  \item 全球核能共识下，中国模式成为可复制的蓝本，能源基础设施竞争格局重塑。
\end{itemize}
\section{宏观与利率}
\textbf{美国通胀的结构性扩散与美联储鹰派预期升温，同时中国经常账户顺差韧性超预期，人民币资产风险溢价正在重构。}

\begin{itemize}
  \item 美国4月核心CPI超预期，租金上涨暴露劳工统计局数据方法论的滞后性，技术类产品开始涨价，供应链成本压力正在向下游传递，通胀正从临时因素转变为结构性成本传递。
  \item 核心PCE通胀年化同比达到3.3\%，显著高于美联储2\%的目标，市场对2026年降息的预期已经降温，甚至开始定价加息。
  \item 中国一季度经常账户顺差占GDP比例3.8\%，服务贸易逆差从720亿美元收窄至440亿美元，旅游服务逆差收窄成为经常账户的新支撑。
  \item 货物贸易顺差维持在2470亿美元，全年顺差收窄幅度预计仅约2\%，远低于市场此前预期的5-10\%。
  \item 人民币资产的风险溢价正在重构，中美元首峰会释放的贸易与投资实质性松动信号，使全球资本对人民币资产的定价逻辑从地缘政治风险折价向结构性合作折价迁移。
\end{itemize}
\begin{figure}[htbp]
\centering
\includegraphics[width=0.88\linewidth]{figures/fig_010.jpg}
\caption{Figure 1 - Figure 1: Imports growth exceeded exports in Q1 2026...}
\end{figure}
\begin{figure}[htbp]
\centering
\includegraphics[width=0.88\linewidth]{figures/fig_011.jpg}
\caption{Figure 2 - Figure 2: ...driven by imports of AI-related products and upstream industrial inputs Imports increase, Q1 2026 v.s. Q1 2025}
\end{figure}
{\small\color{refgray}\textbf{References:} [R002] 进口激增的真正信号：人民币正在被重新定价; [R005] 中国经常账户顺差反弹：市场低估的不是出口，而是服务逆差的收窄能力; [R008] 市场真正低估的不是通胀本身，而是通胀正从“临时因素”转变为“结构性成本传递”; [R011] 市场真正低估的不是贸易缓和，而是人民币资产的风险溢价重构}

\section{AI/算力与科技}
\textbf{AI需求正从训练转向推理，存储芯片进入三重超级周期，中国ERP厂商的AI原生功能成为收入增长新引擎。}

\begin{itemize}
  \item AI驱动的存储需求正在进入指数级增长阶段，用户数、使用时长、任务复杂度、推理token消耗和代理式AI的乘法效应使未来五年总存储需求可能增长数千倍，而供给端只能增长5-6倍。
  \item 存储芯片行业正经历从周期性向成长性的根本转变，三星电子和SK海力士的远期市盈率仍在6倍左右，而台积电约为20倍，估值折价正在失去合理性。
  \item 中国ERP厂商的AI原生功能正在成为拉动用户支出和增量收入的新引擎，金蝶AI原生产品合同达2.3亿元，预计2026年AI套件收入将超过10亿元。
  \item AI对内部研发效率的提升正在改善人均产出，金蝶目标2026年70\%的新代码由AI生成，研发效率有望翻倍。
  \item 中国进口激增中AI相关产品（芯片、半导体设备、计算机）贡献约40\%的增量，表明中国正在深度嵌入全球AI产业链的产能扩张周期。
\end{itemize}
\begin{figure}[htbp]
\centering
\includegraphics[width=0.88\linewidth]{figures/fig_012.jpg}
\caption{Exhibit 2 - Exhibit 2: Kingdee DCF valuation Exhibit 3: Kingdee 12M forward P/S}
\end{figure}
\begin{figure}[htbp]
\centering
\includegraphics[width=0.88\linewidth]{figures/fig_013.jpg}
\caption{Exhibit 5 - Exhibit 5: Yonyou DCF valuation Exhibit 6: Yonyou 12M forward P/S}
\end{figure}
{\small\color{refgray}\textbf{References:} [R002] 进口激增的真正信号：人民币正在被重新定价; [R006] 市场真正低估的不是需求，而是供给侧的再定价; [R009] 中国ERP厂商的真正拐点，不是云化，而是AI对定价权的重塑}

\section{能源与大宗}
\textbf{全球核能共识下中国模式成为蓝本，铝市场出口回暖但内需分化，供给侧刚性约束放大价格含义。}

\begin{itemize}
  \item 35个新国家正在规划核电站，中国在建+规划规模最大（30GW在建+42GW规划），中国凭借低成本、短工期、规模化交付记录成为全球核能扩张的默认参考系。
  \item 俄罗斯在核能出口上领先，已与29个国家有核能合作协议，但中国模式的可复制性正在吸引更多新进入者。
  \item 铝市场出口订单显著修复，4月环比增长约30\%，5月订单量相较3月增长50\%，出口利润每吨提升约1800元人民币。
  \item 国内铝需求呈现结构性分化，龙头加工企业产能利用率保持高位，而中小企业在建筑领域的订单下滑明显，行业洗牌加速。
  \item 电解铝产能的违规超产空间几乎被技术手段封死，供给侧的刚性约束将放大需求再分配的价格含义。
\end{itemize}
\begin{figure}[htbp]
\centering
\includegraphics[width=0.88\linewidth]{figures/fig_001.jpg}
\caption{EXHIBIT 1 - EXHIBIT 1: France and Japan's pivot to nuclear power post 1973 energy crisis France: Historical electricity generation (TWh)}
\end{figure}
\begin{figure}[htbp]
\centering
\includegraphics[width=0.88\linewidth]{figures/fig_002.jpg}
\caption{EXHIBIT 2 - EXHIBIT 2: China accounts for \textbackslash{}\textasciitilde{}50\% of total under construction and planned nuclear plants World Nuclear Capacity (GW)}
\end{figure}
\begin{figure}[htbp]
\centering
\includegraphics[width=0.88\linewidth]{figures/fig_003.jpg}
\caption{EXHIBIT 4 - EXHIBIT 4: Russia takes the lead on partnership for new nuclear capacity addition in 29 countries \# of new countries being supported by existing nuclear majors}
\end{figure}
\begin{figure}[htbp]
\centering
\includegraphics[width=0.88\linewidth]{figures/fig_008.jpg}
\caption{EXHIBIT 9 - EXHIBIT 9: Execution time: Most countries globally take 10 years, China does it in 5 years}
\end{figure}
\begin{figure}[htbp]
\centering
\includegraphics[width=0.88\linewidth]{figures/fig_009.jpg}
\caption{EXHIBIT 10 - EXHIBIT 10: China has one of the lowest LCOEs for nuclear power - given low capex and construction time Nuclear power LCOE (\textbackslash{}\$/MWh)}
\end{figure}
{\small\color{refgray}\textbf{References:} [R001] 全球核能共识的真正赢家，不是技术，而是中国模式; [R010] 铝市场正在经历的不是需求崩塌，而是结构性的再分配}

\section{地产与消费}
\textbf{中国内需呈现延迟而非消失的特征，铝下游在特定价格区间提货加快，但地产链疲弱仍在影响中小企业。}

\begin{itemize}
  \item 铝加工企业反馈，铝价跌至23500-24300元/吨区间时，下游提货明显加快，说明下游不是没需求，而是在等待价格。
  \item 龙头铝板带企业产能提升20-30\%，订单依然饱满，但高度依赖地产的中小加工厂5月需求明显走弱。
  \item 中国ERP厂商的云订阅收入增长故事已经不够，AI才是2026年的增量来源，但企业IT支出仍面临短期压力。
  \item 金蝶2026年第一季度ARR同比增长19\%至42亿元人民币，用友收入同比增长6\%，云收入增长17\%，毛利率提升5.7个百分点。
\end{itemize}
{\small\color{refgray}\textbf{References:} [R009] 中国ERP厂商的真正拐点，不是云化，而是AI对定价权的重塑; [R010] 铝市场正在经历的不是需求崩塌，而是结构性的再分配}

\section{区域市场与外汇}
\textbf{新兴市场外汇面临结构性压力，印尼盾脆弱性突出，而人民币在贸易与投资实质性松动下获得支撑。}

\begin{itemize}
  \item 印尼盾面临外汇储备持续消耗、市场对央行加息意愿存疑、资本外流叠加的三重压力，MSCI指数调整预计带来约17亿美元被动资金流出。
  \item 新加坡元被看好，能源价格推升通胀，3月核心CPI已明显反弹，预计Q3进一步走高。
  \item 美元短期仍有支撑，美国经济韧性仍在，美股创新高吸引资金持续流入美元资产。
  \item 人民币短期有支撑，做空美元/人民币目标价看向6.60，预期收益约3\%，时间窗口到2026年8月底。
  \item 中美元首峰会释放贸易与投资实质性松动信号，包括中国订购200架波音飞机、美国批准H200芯片向中国企业销售、考虑取消约300亿美元中国商品关税等。
\end{itemize}
{\small\color{refgray}\textbf{References:} [R007] 市场真正低估的不是地缘风险，而是新兴市场外汇的结构性压力; [R011] 市场真正低估的不是贸易缓和，而是人民币资产的风险溢价重构}

\section{风险偏好与资产配置}
\textbf{美股盈利增长正从Mag 7向S\&P 493扩散，市场表现将从集中走向分散，周期性板块优于防御性板块。}

\begin{itemize}
  \item S\&P 500未来12个月目标价8300点，基于2027年中滚动EPS达到404美元、估值20.5倍。
  \item 2026年Mag 7净利润增速预计为33\%，而S\&P 493增速将跃升至14\%，盈利增长的中心正在扩散。
  \item 中位数股票的盈利增速创近4年新高，说明上涨不再只靠头部公司，盈利广度正在改善。
  \item 周期性板块优于防御性板块，工业、金融、可选消费被看好，大型云计算公司被视为有吸引力的相对价值交易。
  \item 市场对AI主题的过度集中正在让位于对盈利复苏广度的重新定价，投资者需要关注盈利修正广度等领先指标。
\end{itemize}
{\small\color{refgray}\textbf{References:} [R004] 市场真正低估的不是AI，而是盈利复苏的广度}

\section*{结语}
当前市场主线清晰：AI驱动的盈利扩散与供应链重构正在重塑全球资产定价框架，投资者需从地缘风险折价转向结构性再定价逻辑。
\vfill
{\small\color{refgray}Personal reading notes and learning share only. Not investment advice.}
\end{document}
